%==============================================================
%  lecons/algebre/polynomes_image.tex — Polynômes surjectifs sur C, R, Q
%==============================================================

\lecontitle{Polynômes surjectifs sur \texorpdfstring{$\mathbb{C}$, $\mathbb{R}$, $\mathbb{Q}$}{C, R, Q}}{Algèbre — Un exercice d'oral (Mines-Ponts / ENS, MP)}

%=============================================================
%  § 1 — L'ÉNONCÉ
%=============================================================
\secnum{navyblue}{1}{L'énoncé}

\begin{tcolorbox}[thmbox]
  \textbf{Problème.}\enspace%
  \index{polynôme!surjectif}\index{image d'un polynôme}
  Déterminer les polynômes $P \in \mathbb{C}[X]$ tels que :
  \[
    \text{(a)}\ \ P(\mathbb{C}) = \mathbb{C}, \qquad
    \text{(b)}\ \ P(\mathbb{R}) = \mathbb{R}, \qquad
    \text{(c)}\ \ P(\mathbb{Q}) = \mathbb{Q},
  \]
  où l'on note $P(A) = \{\,P(z) : z \in A\,\}$ l'image directe de $A$ par $P$.
\end{tcolorbox}

\begin{significationintuitive}
La même fonction polynomiale « remplit » ou non son corps d'arrivée selon le
corps que l'on regarde. Sur $\mathbb{C}$, la clôture algébrique rend tout
polynôme non constant surjectif. Sur $\mathbb{R}$, c'est une affaire de
\emph{parité} du degré (les branches partent-elles vers $+\infty$ des deux
côtés ?). Sur $\mathbb{Q}$, l'arithmétique des dénominateurs est bien plus
contraignante : seules les fonctions affines y parviennent.
\end{significationintuitive}

\decsep{}

%=============================================================
%  § 2 — UN OUTIL : LES COEFFICIENTS PAR INTERPOLATION
%=============================================================
\secnum{navyblue}{2}{Un lemme d'interpolation}

\begin{tcolorbox}[thmbox]
  \textbf{Lemme (descente des coefficients).}\enspace%
  \index{interpolation de Lagrange}
  Soit $\mathbb{K} \in \{\mathbb{R}, \mathbb{Q}\}$ et $P \in \mathbb{C}[X]$ de
  degré $\leq n$. Si $P(k) \in \mathbb{K}$ pour tout $k \in \{0,1,\ldots,n\}$,
  alors $P \in \mathbb{K}[X]$.
\end{tcolorbox}

\begin{ideepreuve}
Un polynôme de degré $\leq n$ est entièrement déterminé par ses valeurs en
$n+1$ points : la formule d'interpolation de Lagrange le reconstitue à partir
de $P(0),\ldots,P(n)$ et de polynômes de base à coefficients rationnels.
\end{ideepreuve}

\noindent{\bfseries\color{navyblue} Preuve.}

\begin{mdframed}[style=proofbar]
  Notons $L_0,\ldots,L_n$ les polynômes de Lagrange associés aux points
  $0,1,\ldots,n$ :
  \[
    L_k(X) = \prod_{\substack{0 \leq j \leq n \\ j \neq k}}
             \frac{X - j}{k - j} \in \mathbb{Q}[X].
  \]
  Le polynôme $\widetilde P = \sum_{k=0}^{n} P(k)\,L_k$ vérifie
  $\widetilde P(i) = P(i)$ pour tout $i \in \{0,\ldots,n\}$. Comme $P$ et
  $\widetilde P$ sont de degré $\leq n$ et coïncident en $n+1$ points, ils sont
  égaux. Or chaque $L_k \in \mathbb{Q}[X] \subseteq \mathbb{K}[X]$ et chaque
  $P(k) \in \mathbb{K}$ : donc $P = \widetilde P \in \mathbb{K}[X]$.
  \hfill$\blacksquare$
\end{mdframed}

\decsep{}

%=============================================================
%  § 3 — SUR C : LA CLÔTURE ALGÉBRIQUE SUFFIT
%=============================================================
\secnum{navyblue}{3}{Sur \texorpdfstring{$\mathbb{C}$}{C}}

\begin{tcolorbox}[thmbox]
  \textbf{Réponse (a).}\enspace%
  \index{d'Alembert-Gauss (théorème de)}
  $P(\mathbb{C}) = \mathbb{C}$ si et seulement si $P$ est \emph{non constant}.
\end{tcolorbox}

\noindent{\bfseries\color{navyblue} Preuve.}

\begin{mdframed}[style=proofbar]
  Si $P = c$ est constant, alors $P(\mathbb{C}) = \{c\} \neq \mathbb{C}$.

  \smallskip\noindent
  Réciproquement, supposons $P$ non constant et soit $w \in \mathbb{C}$. Le
  polynôme $P - w$ est non constant, donc, $\mathbb{C}$ étant algébriquement
  clos (théorème de d'Alembert--Gauss), il admet une racine $z \in \mathbb{C}$ :
  $P(z) = w$. Ainsi tout $w$ est atteint, $P(\mathbb{C}) = \mathbb{C}$.
  \hfill$\blacksquare$
\end{mdframed}

\decsep{}

%=============================================================
%  § 4 — SUR R : UNE QUESTION DE PARITÉ
%=============================================================
\secnum{navyblue}{4}{Sur \texorpdfstring{$\mathbb{R}$}{R}}

\begin{tcolorbox}[thmbox]
  \textbf{Réponse (b).}\enspace%
  $P(\mathbb{R}) = \mathbb{R}$ si et seulement si $P \in \mathbb{R}[X]$ et
  $\deg P$ est \emph{impair}.
\end{tcolorbox}

\noindent{\bfseries\color{navyblue} Preuve.}

\begin{mdframed}[style=proofbar]
  \textit{Coefficients réels.}\enspace
  Si $P(\mathbb{R}) = \mathbb{R}$, alors en particulier $P(k) \in \mathbb{R}$
  pour $k \in \{0,\ldots,\deg P\}$, donc $P \in \mathbb{R}[X]$ par le lemme du
  § 2. On peut désormais voir $P$ comme une fonction $\mathbb{R} \to \mathbb{R}$,
  continue.

  \smallskip\noindent
  \textit{Le degré pair échoue.}\enspace
  Si $\deg P = 2p$ est pair, de coefficient dominant $a$, alors
  $P(x) \sim a\,x^{2p}$ quand $x \to \pm\infty$ : $P$ tend vers $+\infty$ (si
  $a>0$) ou $-\infty$ (si $a<0$) aux deux extrémités. Étant continue, $P$ est
  alors minorée (resp. majorée) sur $\mathbb{R}$, donc \emph{non surjective}.
  Le cas constant ($\deg P = 0$) est exclu de même.

  \smallskip\noindent
  \textit{Le degré impair convient.}\enspace
  Si $\deg P$ est impair, $P$ admet des limites infinies de signes opposés en
  $-\infty$ et $+\infty$. Par le théorème des valeurs intermédiaires, $P$ prend
  toute valeur réelle : $P(\mathbb{R}) = \mathbb{R}$.
  \hfill$\blacksquare$
\end{mdframed}

\decsep{}

%=============================================================
%  § 5 — SUR Q : SEULES LES FONCTIONS AFFINES
%=============================================================
\secnum{navyblue}{5}{Sur \texorpdfstring{$\mathbb{Q}$}{Q}}

\begin{tcolorbox}[thmbox]
  \textbf{Réponse (c).}\enspace%
  \index{théorème des racines rationnelles}
  $P(\mathbb{Q}) = \mathbb{Q}$ si et seulement si $P = aX + b$ avec
  $a \in \mathbb{Q}^*$ et $b \in \mathbb{Q}$.
\end{tcolorbox}

\begin{significationintuitive}
Sur $\mathbb{Q}$, le critère de parité ne suffit plus : $X^3$ est de degré
impair mais $2$ n'est pas un cube de rationnel, donc $X^3$ n'atteint pas $2$.
La raison profonde est arithmétique — les dénominateurs des antécédents sont
bornés par le coefficient dominant, ce qui empêche dès le degré $2$ d'atteindre
les rationnels de dénominateur arbitrairement grand.
\end{significationintuitive}

\noindent{\bfseries\color{navyblue} Preuve.}

\begin{mdframed}[style=proofbar]
  \textit{Coefficients rationnels et cas affine.}\enspace
  Si $P(\mathbb{Q}) = \mathbb{Q}$, alors $P(k) \in \mathbb{Q}$ pour
  $k \in \{0,\ldots,\deg P\}$, donc $P \in \mathbb{Q}[X]$ (lemme du § 2).
  Réciproquement, si $P = aX + b$ avec $a \neq 0$ rationnels, tout
  $y \in \mathbb{Q}$ admet l'antécédent rationnel $x = (y-b)/a$ : $P$ est même
  une bijection de $\mathbb{Q}$. Il reste à exclure $\deg P \geq 2$.

  \smallskip\noindent
  \textit{Une borne sur les dénominateurs des antécédents.}\enspace
  Supposons $\deg P = n \geq 2$ et écrivons
  $P = \frac{1}{d}\bigl(c_n X^n + \cdots + c_1 X + c_0\bigr)$ avec
  $c_i \in \mathbb{Z}$, $c_n \neq 0$, $d \in \mathbb{Z}_{>0}$. Soit $y \in P(\mathbb{Q})$
  et $x = p/m$ un antécédent écrit sous forme irréductible ($m \geq 1$,
  $\pgcd(p,m)=1$). En multipliant $P(p/m) = y$ par $d\,m^n$ :
  \[
    c_n p^n + c_{n-1}p^{n-1}m + \cdots + c_0 m^n = d\,y\,m^n .
  \]
  Tous les termes sauf $c_n p^n$ sont divisibles par $m$, ainsi que le membre de
  droite ($m \mid m^n$ car $n \geq 1$). Donc $m \mid c_n p^n$, et comme
  $\pgcd(p,m)=1$, il vient $m \mid c_n$.

  \smallskip\noindent
  \textit{L'image rate des rationnels.}\enspace
  Les antécédents possibles ont donc un dénominateur $m$ divisant l'entier fixe
  $c_n$ : ils vivent dans $\frac{1}{c_n}\mathbb{Z}$. Pour un tel $x = t/c_n$
  ($t \in \mathbb{Z}$),
  \[
    y = P(x) = \frac{1}{d}\sum_{i=0}^{n} c_i \frac{t^i}{c_n^{\,i}}
      \in \frac{1}{d\,c_n^{\,n}}\,\mathbb{Z}.
  \]
  Autrement dit $P(\mathbb{Q}) \subseteq \frac{1}{d\,c_n^{\,n}}\mathbb{Z}$, un
  sous-groupe \emph{discret} de $\mathbb{Q}$. Or $\frac{1}{2\,d\,c_n^{\,n}} \in \mathbb{Q}$
  n'y appartient pas : $P$ n'est pas surjective. Donc $\deg P \leq 1$, ce qui
  achève la caractérisation. \hfill$\blacksquare$
\end{mdframed}

\rubriq{Exemples et contre-exemples}

\begin{mdframed}[style=exbar]
  \textbf{1. $P = X^2$ sur $\mathbb{C}$ et $\mathbb{R}$.}\enspace%
  Surjectif sur $\mathbb{C}$ (non constant), mais $P(\mathbb{R}) = \mathbb{R}_+ \neq \mathbb{R}$
  (degré pair).

  \smallskip\noindent
  \textbf{2. $P = X^3 - X$.}\enspace%
  Degré impair à coefficients réels : $P(\mathbb{R}) = \mathbb{R}$.

  \smallskip\noindent
  \textbf{3. $P = 2X + 1$.}\enspace%
  Convient sur les trois corps : $P(\mathbb{C})=\mathbb{C}$, $P(\mathbb{R})=\mathbb{R}$
  et $P(\mathbb{Q})=\mathbb{Q}$.
\end{mdframed}

\noindent{\bfseries\color{counterred} Contre-exemples instructifs.}

\begin{mdframed}[style=cexbar]
  \textbf{Le degré impair ne suffit pas sur $\mathbb{Q}$.}\enspace%
  $P = X^3$ vérifie $P(\mathbb{Q}) \subseteq \mathbb{Q}$ et $\deg P$ impair,
  mais $2 \notin P(\mathbb{Q})$ : $\sqrt[3]{2}$ est irrationnel. La parité,
  décisive sur $\mathbb{R}$, est insuffisante sur $\mathbb{Q}$.

  \smallskip\noindent
  \textbf{Coefficients complexes.}\enspace%
  $P = iX$ vérifie $P(\mathbb{C}) = \mathbb{C}$, mais $P(\mathbb{R}) = i\mathbb{R} \neq \mathbb{R}$
  et $P(\mathbb{Q}) = i\mathbb{Q} \neq \mathbb{Q}$ : sans coefficients réels
  (resp. rationnels), l'image quitte la droite réelle.
\end{mdframed}

\exolabel

\begin{mdframed}[style=exobar]
  \begin{enumerate}
    \item Montrer que $P \in \mathbb{C}[X]$ vérifie $P(\mathbb{R}) \subseteq \mathbb{R}$
          si et seulement si $P \in \mathbb{R}[X]$.
    \item Déterminer les $P \in \mathbb{C}[X]$ tels que $P(\mathbb{R}) = \mathbb{R}_+$.
    \item Soit $P \in \mathbb{Q}[X]$ de degré $\geq 2$. Montrer que
          $\mathbb{Q} \setminus P(\mathbb{Q})$ est infini.
    \item Caractériser les $P \in \mathbb{C}[X]$ tels que $P(\mathbb{Z}) \subseteq \mathbb{Z}$
          (polynômes à valeurs entières~: base des $\binom{X}{k}$).
    \item Existe-t-il $P \in \mathbb{C}[X]$ non constant tel que
          $P(\mathbb{U}) = \mathbb{U}$, où $\mathbb{U}$ est le cercle unité ?
          (Indication~: $P = X^n$.)
  \end{enumerate}
\end{mdframed}
